\section{The degree and the relation ideal}\label{sec:main-statement}
Put $r=\rho+\tau+3\#\mathcal C$, as in Theorem~\ref{thm:main}.
A transcendence basis contains a first-level residue basis.
It also contains all distinct levels $(\eta,k)$ with $k\ge2$.
For each class it contains one triple:
$(I_c,V_c,S_c)$, $(F_{s,a},F'_{s,a},\Sigma_{s,a})$, or
$(f_{\nu,c},f'_{\nu,c},f_{-\nu,c})$. It stays independent over $T(L)$
when $e_{(0,0)}\notin\mathcal V_{\rm res}$. Then $L$ is transcendental over
$K_{\rm full}$. Otherwise $L\in K_{\rm full}$.

Use the coordinate variables defined before Theorem~\ref{thm:main}.
The full prime ideal $\mathfrak p\subset T[X]$ has dimension $r$.
Its generators are:
\begin{enumerate}\setlength{\itemsep}{0pt}\setlength{\parsep}{0pt}
\item $X_e-t^{v_e}$ for listed exponentials of rate $\sum(v_e)_h\lambda_h$;
\item for nonbasis first-level $F_j$,
\begin{equation}\label{eq:residue-relations}
F_j-\sum_i\kappa_{ji}F_i-r_j-c_j,\quad
\operatorname{res}(F_j)=\sum_i\kappa_{ji}\operatorname{res}(F_i),\quad r_j\in\C(z)[t^{\pm1}],\ c_j\in\C;
\end{equation}
\item one representative determinant,
\begin{equation}\label{eq:determinant-relations}
\begin{array}{ll}
\operatorname{Bess}(c,0):&z(IW-VS)-I^2+1,\\
\operatorname{Kum}(s,a):&z(F\Sigma'-F'\Sigma)-F^2+e^{sz},\\
\operatorname{Bess}(c,\nu):&z(f_\nu f'_{-\nu}-f'_\nu f_{-\nu})-2\nu f_\nu f_{-\nu}+2\nu;
\end{array}
\end{equation}
\item every nonrepresentative coordinate minus its $T$-linear representative
expression from Theorem~\ref{thm:linkage} and derivatives. The last determinant
uses its own nonresonant representative order. \end{enumerate}
Include $X_q-X_{\operatorname{rep}(q)}$ for every duplicate slot of any coordinate type.
For each listed constant, include $X_c-c$.
These monic graphs eliminate duplicate and constant slots.
The degree is unchanged. Only distinct higher levels remain free.
The data definitions preceding Theorem~\ref{thm:main} define $\operatorname{rep}$.
For half-integer orders, add the exponential graphs ($\pm c\in\Gamma$)
\[
 X_{f_{1/2,c}}-\frac{\sinh(cz)}{cz},\qquad
 X_{f_{-1/2,c}}-\cosh(cz),
\]
and add their derivatives and order shifts from \eqref{eq:link-order}.
Write $\sinh,\cosh$ as Laurent polynomials in the basis exponentials.

For the Laurent form, replace $t$ by $\Theta$ in all generators.
Work over $\mathscr A_{\exp}=\C(z)[\Theta^{\pm1}]$.
The generated prime ideal is exactly
$\mathfrak P=\ker(\mathscr A_{\exp}[X]\to K_{\rm full})$ under $\Theta_h\mapsto t_h$,
of dimension $d+r$. \mbox{Its} contraction to $\C(z)[X]$ is the full coordinate ideal.
Laurent saturation and elimination compute it, including exponential toric relations.
To replace $\Theta$-monomials by listed exponentials, first clear negative powers.
Then saturate by their product. Let $k_1$ be any field containing
the data, listed constants and fundamental $c_j$.
The same generators give $\mathfrak P\cap k_1(z)[\Theta^{\pm1}][X]$.
This ideal is prime of dimension $d+r$.
The corresponding degree is $\trdeg_{k_1(z)}k_1(z)(t,\mathcal F)=d+r$.
For the normalization-compatible input of Appendix~\ref{app:algorithm},
the fundamental constants $c_j$ belong to $k_1$.
Under Theorem~\ref{thm:values}, first-level coordinates vanish at zero.
The generators lie in
$\Qbar[z^{\pm1}][\Theta^{\pm1}][X]$.

